\section{Which signatures a reduction must retain}\label{sec:signatures}
Every exponent greater than two has a useful small divisor. The divisor need not be the same for the three exponents.
\begin{proposition}[Prime or four reduction]\label{prop:primefour}
If a primitive solution of \eqref{eq:beal} exists, there is a primitive solution whose three exponents each belong to
\[
\{4\}\cup\{p:p\text{ is an odd prime}\}.
\]
Conversely, absence of primitive solutions for every ordered signature in this set implies Beal.
\end{proposition}
\begin{proof}
Let $n>2$. If $n$ has an odd prime divisor $p$, write $D^n=(D^{n/p})^p$. If it has no odd prime divisor, it is a power of two at least four, so write $D^n=(D^{n/4})^4$. Apply this operation to each term separately. The new bases have exactly the same prime supports as the old bases, so primitivity is preserved. The converse follows by reducing any hypothetical primitive solution in this way.
\end{proof}
The proposition retains infinitely many signatures. It does not reduce them all to $(3,5,7)$. A theorem for that signature also handles signatures $(3a,5b,7c)$ after absorbing the multipliers into the bases, and the symmetric version with the first two terms exchanged. It does not handle an unrelated prime signature such as $(5,7,11)$ merely because the same packet language can describe both.

If all three exponents have a common divisor $n>2$, absorbing the remaining factors produces an equal exponent Fermat equation. Fermat's Last Theorem then excludes the primitive solution \cite{wiles1995,taylor-wiles1995}. This use requires a common divisor of all three exponents. Sharing a divisor in only one pair is weaker. The general prime or four reduction must therefore retain mixed signatures even after the equal exponent cases are removed.

The fixed signature finiteness theorem of Darmon and Granville provides another useful perspective \cite{darmon-granville1995}. In the hyperbolic range
\[
\frac1x+\frac1y+\frac1z<1,
\]
their theorem gives finiteness of primitive solutions for a fixed equation and fixed exponents. For $x,y,z>2$, the only nonhyperbolic possibility is $(3,3,3)$, where the reciprocal sum is one and Fermat already applies. Finiteness is valuable, but it is not absence. It also does not give a uniform finite enumeration over all signatures.

\subsection{The conditional abc height bound}
An abc estimate makes the role of the signature explicit. Let
\[
H=C^z,\qquad R=\rad(ABC),\qquad s=\frac1x+\frac1y+\frac1z.
\]
For a primitive solution, the powered terms are pairwise coprime and their radical is $R$. Positivity gives $ABC<H^s$. Assume the all epsilon abc estimate
\begin{equation}\label{eq:abc-assumed}
H\le K_\eps R^{1+\eps}
\end{equation}
for every $\eps>0$, with one constant $K_\eps$ independent of the triple. Then
\[
H<K_\eps H^{(1+\eps)s}.
\]
Whenever $(1+\eps)s<1$, it follows that
\begin{equation}\label{eq:signatureheight}
H<K_\eps^{1/(1-(1+\eps)s)}.
\end{equation}
This is a conditional bound, not an unconditional result obtained from the exponent equation alone.

For $(3,5,7)$, the reciprocal sum is $71/105$. Thus $0<\eps<34/71$ gives
\begin{equation}\label{eq:357abc}
H<K_\eps^{105/(34-71\eps)},\qquad
C<K_\eps^{15/(34-71\eps)}.
\end{equation}
At $\eps=1/10$, the latter exponent is $150/269$. To turn this into a proof of absence would require the abc input, an explicit usable constant, and a complete verification below the resulting bound.

There is a stronger uniform bound across the Beal range. Apart from $(3,3,3)$, already excluded by Fermat, one exponent is at least four. Thus $s\le1/3+1/3+1/4=11/12$. Take $\eps=1/22$ and enlarge its constant to $K_{1/22}\ge1$. Then $(1+\eps)s\le23/24$, and \eqref{eq:signatureheight} gives
\begin{equation}\label{eq:uniformbealabc}
H<K_{1/22}^{24}.
\end{equation}
This conditional constant is independent of the signature. It bounds the powered terms and yields finitely many possible primitive powered triples. It does not show that the finite set is empty. Since $2^z\le H$, it also bounds $z$ and excludes solutions once the minimum exponent is sufficiently large. The all epsilon hypothesis and the gap between finiteness and absence remain essential \cite{darmon-granville1995}.

Unit bases can be disposed of separately by the classical theorem on consecutive perfect powers. If $A=1$ and $B>1$, then $C^z-B^y=1$. Mih\u{a}ilescu's theorem gives the only pair $3^2-2^3=1$, whose upper exponent is two and violates the Beal range \cite{mihailescu2004,mihailescu2005,bilu2005}. The same argument handles $B=1$. If both are one, $C^z=2$ is impossible. Hence every actual Beal solution has all bases greater than one. After this accepted deep input, a powered height bound also bounds every exponent label. This is a use of a classical theorem, not a new proof of it.

\subsection{Why a generic radical ceiling fails}
The exponent equation gives $R<H^{71/105}$ for the signature $(3,5,7)$. It is tempting to reverse this relation by postulating $H\le R^{105/71}$ for every primitive sum. That universal assertion is false. The exact equality
\[
2+109\cdot3^{10}=23^5=6\,436\,343
\]
has radical $2\cdot3\cdot109\cdot23=15\,042$, and
\[
(23^5)^{71}>15\,042^{105}.
\]
Its logarithmic quality is about $1.629911684$. It is not a $(3,5,7)$ counterexample. It shows that an estimate strong enough to close the desired signature must retain its power conditions or contain the constants and exceptions that a genuine abc type estimate allows.
